\newcommand{\C}{\mathbb{C}}
\newcommand{\R}{\mathbb{R}}
\newcommand{\Z}{\mathbb{Z}}
\newcommand{\Tr}{\operatorname{Tr}}
\newcommand{\Impart}{\operatorname{Im}}
\newcommand{\AF}{\operatorname{AF}}
\newcommand{\AX}{\operatorname{AX}}
\newcommand{\GL}{\operatorname{GL}}

\examyear{2012}
% \examera{平成24年}
% \examfield{解析}
% \examgroup{B}
% \examnumber{12}
% \examtitle{平成24年 B 第12問}
% \examtag{Fourier解析}
% \examtag{周期関数}
% \examtag{級数}

\(f\) は \([0,+\infty)\) 上の実数値連続関数であり，
\[
  f(0)\ne0,\qquad f(x+1)=f(x)\quad (x\in[0,+\infty))
\]
をみたす。\(t>0\) に対して
\[
  F(t)=\sum_{n=-\infty}^{+\infty}
  \left|\int_0^t f(x)e^{2\pi i n x}\,dx\right|^2
\]
と定義する。ただし \(i=\sqrt{-1}\) である。

\begin{enumerate}
  \item \(F(t)<+\infty\ (0<t<+\infty)\) であることを証明せよ。
  \item \(F(t)\) が微分可能となる点 \(t>0\) をすべて求め，その点 \(t\) における \(F'(t)\) を求めよ。
\end{enumerate}

%